\renewcommand*{\authorOfNextLines}{Helmut Quinz}
Die Gesamtbeurteilung der Diplomarbeit setzt sich aus den Bereichen
\begin{itemize}
	\item praktische Arbeit und schriftliche Dokumentation sowie
	\item Präsentation und Verteidigung
\end{itemize}
zusammen.\\
Für jeden Bereich sind spezifische Kriterien definiert die vom Betreuer als
\begin{description}
	\item[4] überwiegend erfüllt,
	\item[3] zur Gänze erfüllt,
	\item[2] darüber hinausgehend erfüllt und
	\item[1] weit darüber hinausgehend erfüllt 
\end{description}
eingeschätzt werden.\\
Diese Einschätzungen geben dem Betreuer eine Grundlage für die Bewertung der Diplomarbeit in einer Ziffernnote zwischen 1 (Sehr gut) und 5 (Nicht genügend).
\paragraph{Hinweis:}~\\
Auch wenn zwischen Abgabe der Dokumentation und Verteidigung ein erheblicher Zeitraum liegt, 
ist es dem Betreuer rechtlich nicht gestattet mit dem Kandidaten über die Beurteilung des praktischen Teiles zu sprechen.